% tab:sources
\begin{table}[htbp]\centering\scriptsize\setlength{\tabcolsep}{3pt}
\caption{Indicators and sources of the Latvian Soviet-period series, 1937--95}
\label{tab:sources}
\begin{threeparttable}
\begin{tabular}{p{2.9cm}p{11.3cm}c}
\toprule
Sector & Indicator and main sources & Appendix \\
\midrule
Crops & Physical output of 9--10 crop categories, 1937--95; 1990 value-added weights from the World Bank's 1993 study, grains split with the 1995 Eurostat agricultural accounts; sown areas for 1945--55 from unpublished returns & A.1 \\
Livestock & Physical output of 7--16 animal products; 1941--9 yields obtained from animal numbers and the yields per animal of 1940, 1945 and 1950 & A.1 \\
Forestry and fishing & Timber cut (cubic metres) and fish catch (tonnes); 1941, 1971 and 1974--6 catches interpolated & A.1 \\
Industry & 129 physical product series in 11 branches (Table~\ref{tab:products}); product weights from the 1987 input--output study and the 1993 World Bank study, products consolidated into 22 minor branches by index averages because no archival late-1980s output valuation survives for Latvia; branch weights from the 1990 benchmark & A.2 \\
Construction & Weighted index of material inputs into construction, 1987 input--output weights, following the CIA procedure & A.3 \\
Transport and communications & Freight tonne-kilometres and passenger-kilometres, all modes; postal, telegraph and telephone traffic and radio and television sets; 1945--50 from the Russian State Archive of the Economy; 1941--4 assumed to move with the rest of the economy; 1990s from sectoral employment & A.4 \\
Trade, accommodation and catering & Soviet retail and wholesale turnover indices at comparable prices (the only value series used), chained across base-year changes; early post-war years from the USSR domestic-trade compendia & A.5 \\
Finance, insurance and housing & Finance and insurance: employment by branch, adjusted for productivity with the ratio of trade turnover to trade employment, as the CIA did; housing: housing stock derived from construction and population & A.5 \\
Public administration, health, education, science & Employment by branch from Soviet yearbooks, the Ostland statistical bulletins and World Bank studies; 1941--4, 1946--9 and 1951--4 gaps interpolated & A.6 \\
Defence & Soviet garrison of 30,000 in 1940 and 80,000 in 1990 plus the republican forces and police of 1940; 1941--89 interpolated between the 1940 and 1990 totals in step with public-administration employment; 1991--5 from the withdrawal schedule and the new Latvian forces & A.6 \\
Other services & Employment & A.6 \\
Population & Year-end estimates and censuses of 1935, 1959, 1970, 1979 and 1989; wartime and deportation losses from the demographic literature; Soviet garrison added in 1940--1; present-day borders from 1945 & A.7 \\
\bottomrule
\end{tabular}
\begin{tablenotes}[para,flushleft]
\footnotesize
\textit{Notes:} the appendix column refers to the sections of the online appendix in which every series is documented year by year with its sources: the Soviet Latvian statistical yearbooks (\textit{Latvijas PSR tautas saimniec\={\i}ba}, 1957--90), the \textit{Statistische Berichte f\"{u}r das Ostland} (1942--4) and \textit{Ostland in Zahlen} (1942), the USSR domestic-trade compendia (\textit{Sovetskaya torgovlya}, 1964 and 1989), the World Bank's \textit{Latvia: the transition to a market economy} (1993) and \textit{Statistical handbook: states of the former USSR} (1994--6), the Russian State Archive of the Economy (RGAE) and the Latvian State Archives. \textit{Sources:} see online appendix.
\end{tablenotes}
\end{threeparttable}
\end{table}
